\chapter{Background and Related Work}\label{ch:related}

\section{Computational models of expressive performance}
Expressive-performance research models how performers deviate from the
notated score in timing, dynamics and articulation. Historically this was
done with rule-based systems, and more recently with sequence models trained
on aligned corpora \cite{cancino2018,oore2018}. Most such work predicts a
single expressive rendering or studies aggregate statistics. This thesis
differs in framing the problem as per-note Bayesian inference with
calibrated uncertainty. It asks not only how accurately expression can be
predicted but also how \emph{trustworthy} the per-note confidence is.

\section{Score-to-performance alignment}
Aligned corpora pair a symbolic score with a performance through a note- or
beat-level correspondence, from which continuous deviations can be read off.
ASAP provides score$\leftrightarrow$performance alignments for piano with
beat and downbeat annotations \cite{asap}, and it is the source of the
targets used here. The quality of this alignment bounds from above the
quality of the extracted targets. This dependence is treated as an explicit
limitation of the work (Chapter~\ref{ch:discussion}).

\section{Structured priors: graph Gaussian processes, GMRFs, and the SPDE approach}
Gaussian Markov random fields place a Gaussian prior whose sparse precision
encodes conditional-independence structure \cite{rueheld2005}. The SPDE
approach links Gaussian fields to GMRFs and yields Mat\'ern-type priors with
interpretable range and smoothness \cite{lindgren2011}. These constructions
extend to graphs, giving Mat\'ern Gaussian processes on graph nodes
\cite{borovitskiy2021}. For vector-valued outputs, both the
intrinsic-coregionalization construction of \cite{borovitskiy2021} and the
Kronecker output coupling of Gaussian processes over graphs
\cite{venkitaraman2018} prescribe an explicit cross-output covariance.
Gaussian processes on \emph{evolving} graphs \cite{blancomulero2021} address
dynamic structure, which a fixed score does not need. The proposed model is
exactly such a multi-output graph Gaussian process, with the graph built
from musical adjacency rather than physical space. Its graph kernel is a
spectral function of the score-graph Laplacian, its channels are coupled by
a coregionalization matrix, its side information is marginalized into the
kernel as Bayesian linear means, and every hyperparameter is learned by the
exact evidence.

\section{Symbolic music language models}
Transformer and RNN models over tokenised MIDI, among them the Music
Transformer, REMI and performance RNNs, learn strong distributions over
symbolic music \cite{musictransformer,remi,oore2018}. This thesis builds a
small decoder-only model from scratch, in the spirit of nanoGPT
\cite{nanogpt}, rather than using a pretrained one. The reason is that the
uncertainty layer must attach to the model's internals. With a model built
here, the calibration of its outputs and the absence of leakage in its
read-out remain inspectable.

\section{Differentiable synthesis}
Differentiable digital signal processing renders audio from interpretable
controls, namely fundamental frequency, harmonic amplitudes and noise,
through a differentiable synthesizer \cite{ddsp}. This is the intended
waveform likelihood of the thesis. It connects the latent performance
variables to a waveform while keeping the amplitudes conditionally
Gaussian, so that they can be marginalised in closed form
(Chapter~\ref{app:phase3}). The DDSP line has since grown hierarchical
note-to-performance-to-audio \emph{generation} (MIDI-DDSP
\cite{wu2022midi}), physically structured interpretable effect modules such
as a differentiable guitar amplifier \cite{yeh2024ddspamp}, and semantically
disentangled timbre control for multitimbral synthesis \cite{dalri2026}.
Those works run the synthesis direction, with learned neural encoders and
decoders. This thesis runs the \emph{inverse analysis} direction. The
synthesizer is a hand-structured probabilistic observation model, a harmonic
basis with Gaussian amplitudes marginalised exactly, rather than a trained
neural decoder, and the recovered object is the posterior over
score-aligned performance variables rather than audio.

A complementary line puts the periodicity itself into a Gaussian-process
\emph{kernel} rather than into a deterministic basis. Spectral mixture
kernels learn a stationary covariance as a Gaussian mixture over
frequencies \cite{wilson2013}. Their non-stationary generalization, the
\emph{generalized spectral mixture} kernel, lets the frequencies, amplitudes
and coherence lengths vary with the input, each modelled as a latent
function \cite{remes2017}. Applied to music audio, per-note change windows
that switch between Gaussian processes with harmonic, exponentiated-cosine
covariances support pitch estimation and gap imputation directly on the
waveform \cite{alvarado2016}. These families are the natural priors for the
\emph{within-note} structure that this thesis measures but does not model,
such as a vibrato whose period and extent drift. Both places where they
could enter, a curve-level estimator and the pitch-curve prior of the
waveform setting, were tried on the development set, and neither improved
on the simpler alternative (\S\ref{ch:future}), so their imprint on this
thesis is conceptual. The design difference is a choice and is worth
stating. Kernel-side constructions marginalize the note's structure away,
whereas this thesis keeps named per-note variables ($c$, $\gamma$,
$f^{\mathrm{vib}}$, \dots) because the variables themselves are the object
of study.

% =====================================================================
